% Title-block partial for the qkit syllabus template.
% Activated via "template-partials: before-body.tex" in the YAML header.
%
% This replaces Quarto's stock before-body.tex, which does nothing beyond
% \maketitle and an optional abstract. It adds the two things that need
% values out of the YAML header, and so cannot live in preamble.tex:
%
%   1. an optional logo above the title (`logo:`, `logo-width:`), and
%   2. the course-information table (`email:`, `web:`, `officehours:`,
%      `classhours:`, `office:`, `classroom:`).
%
% Every field is optional; anything left out of the YAML renders as TBD,
% so the table keeps its shape while a course is still being scheduled.

$if(logo)$
% titling's \maketitle expands \pretitle and \posttitle around the title,
% which is what puts the logo inside the title block rather than on a page
% of its own. \posttitle restores titling's default trailing \vskip.
\pretitle{%
  \begin{center}%
  \LARGE
  \includegraphics[width=$if(logo-width)$$logo-width$$else$0.30\textwidth$endif$]{$logo$}\\[\bigskipamount]%
}
\posttitle{\par\end{center}\vskip 0.5em}
$endif$

$if(term)$
% `term:` rather than `date:`, because Quarto parses `date:` as a calendar
% date and normalises it -- "Fall 2026" comes back out as 2026-01-01. An
% unrecognised key is passed through untouched, so the term prints as
% written, both here and in the running head. Setting `date:` instead
% still works; Quarto's own \date{} is simply left in place.
\date{$term$}
$endif$

$if(title)$
\maketitle
\thispagestyle{firststyle}
$if(abstract)$
\begin{abstract}
$abstract$
\end{abstract}
$endif$
$endif$

% tabularx with two X columns, NOT tabular* with l/r columns.
%
% The py4Fin original used `tabular*` with \extracolsep{\fill}, which
% cannot wrap: an l column and an r column are set at their natural width,
% so the moment one row's two values are together wider than \textwidth
% the row runs past the right margin. "Office Hours: Tuesdays 14:00--16:00,
% or by appointment" against "Class Hours: Monday and Wednesday,
% 10:30--11:50" overflowed by 54.8pt on the stock values alone -- and the
% only symptom is an Overfull \hbox in the log, since nothing in the PDF
% text says the words are outside the margin.
%
% Two X columns each take half of \textwidth and wrap. \raggedright on the
% left and \raggedleft on the right keeps the py4Fin look -- labels flush
% to the outer margins, whitespace collecting in the middle -- while long
% values now spill onto a second line instead of off the page.
%
% \arraybackslash restores \\ as the row separator, which \raggedright and
% \raggedleft otherwise redefine to their line-break meaning.
%
% `web:` is a BARE host -- the https:// below is supplied here. Writing a
% full URL in the YAML produces https://https://... in the link target.
\noindent\begin{tabularx}{\textwidth}{@{}>{\raggedright\arraybackslash}X >{\raggedleft\arraybackslash}X@{}}
E-mail: $if(email)$\texttt{$email$}$else$TBD$endif$ &
  Web: $if(web)$\href{https://$web$}{\texttt{$web$}}$else$TBD$endif$ \\
Office Hours: $if(officehours)$$officehours$$else$TBD$endif$ &
  Class Hours: $if(classhours)$$classhours$$else$TBD$endif$ \\
Office: $if(office)$$office$$else$TBD$endif$ &
  Classroom: $if(classroom)$$classroom$$else$TBD$endif$ \\
 & \\
\hline
\end{tabularx}

\vspace{2mm}
